\documentclass[10pt,a4paper]{amsart}
\usepackage[margin=19mm]{geometry}
\usepackage{amsmath,amssymb,mathtools}
\usepackage[scr=rsfs,cal=euler]{mathalfa}
\usepackage{xfrac}
\usepackage[unicode,hidelinks,bookmarks=false]{hyperref}
\newcommand{\ie}{i.e.\ }
\newcommand{\cf}{cf.\ }
\newcommand{\resp}{respectively\ }
\newcommand{\Subsection}{\subsection}
\providecommand{\nobreakdash}{\nobreak\hbox{-}}
\setlength{\emergencystretch}{2em}
\multlinegap0pt
\usepackage{bm}
\usepackage{bbm}
\usepackage{dsfont}
\usepackage{xspace}
\usepackage{cancel}
\usepackage{mathtools}
\usepackage{amsthm}
\makeatletter
\@ifundefined{theorem}{\newtheorem{theorem}{Theorem}}{}
\makeatother
\title[Pick-up Sticks and the Fibonacci Factorial]{Pick-up Sticks and the Fibonacci Factorial}
\author{Aidan Sudbury, Arthur Sun, David Treeby, Edward Wang}
\date{Source: arXiv:2504.19911v4 (CC BY 4.0)}
\begin{document}
\maketitle

\noindent\emph{Source and licence.} Pick-up Sticks and the Fibonacci Factorial. Version arXiv:2504.19911v4 (CC BY 4.0), distributed under the Creative Commons Attribution 4.0 International licence, as recorded in the arXiv licence metadata for this exact version and shown on the abstract page https://arxiv.org/abs/2504.19911v4. Full rights notice: licenses/arxiv-2504.19911-CC-BY-4.0.txt.\par\smallskip

\noindent\textit{Editorial scope.} Counted unit: the Theorem whose source label is \texttt{None} in the article (source0.tex lines 149--153, source label \texttt{None}), delivered together with the article's own complete original proof of it (source0.tex lines 154--210, one \texttt{proof} environment of 57 lines). No mathematics is written, repaired or re-derived: statement and proof are the article's own text line for line. Required context retained uncounted: none. Every definition, display and label of those lines is retained unchanged. Omitted from this package and not redistributed: 1--148, 211--227. Nothing inside a retained excerpt is omitted or altered: every retained body is delivered exactly as the article's lines are written, so no omission receipt of any kind accompanies this package. Citations: the retained excerpts carry no citation command at all, so no bibliography entry of the article is transcribed and no reference is redistributed. Reference adaptation: none was needed and none was made; every reference inside the delivered unit resolves inside it, so no unresolved reference or citation is produced. Printed slips kept literal: no printed word, symbol, hypothesis or number of the counted statement or of its proof was altered. Layout and class: the article's own class is \texttt{amsart} with packages including mathtools, amssymb, amsthm, microtype, hyperref; the wrapper is the lane's pinned \texttt{amsart} editorial wrapper at 10pt on A4, which restates the theorem-like environments the retained text needs and adds the article's own macro definitions that the retained text needs (none), with extra wrapper packages (bm, bbm, dsfont, xspace, cancel, mathtools). The delivered numbering is the unit's own. Assets: no retained span contains a figure environment, a graphics-inclusion command, a PDF-page inclusion, a table or a TikZ picture; no image file is copied into this package, the package declares no asset dependency and no image file is redistributed with it. Licence: version arXiv:2504.19911v4 of this article is distributed under the Creative Commons Attribution 4.0 International licence, as recorded in the arXiv licence metadata for this exact version and shown on the abstract page https://arxiv.org/abs/2504.19911v4. A full rights notice is delivered at licenses/arxiv-2504.19911-CC-BY-4.0.txt. Characters: every retained excerpt is pure ASCII, so the delivered engine, XeLaTeX with its default Latin Modern text font, reports no missing character.
\begin{theorem}
  If \(n\) real numbers are chosen uniformly from the unit interval \([0, 1]\), the probability that no four of the numbers can form a quadrilateral is \[
    \frac{1}{T_n - T_{n-2}} \prod_{i = 1}^{n-1} \frac{1}{T_i}
  ,\] where \(T_k\) are the \emph{Tribonacci} numbers defined by \(T_1 = 1\), \(T_2 = 1\), \(T_3 = 2\), \(T_{k} = T_{k-1} + T_{k-2} + T_{k-3}\) for \(k > 3\).
\end{theorem}

\begin{proof}
  Like before, the inequality can be reduced to simultaneously satisfying \[
    U_{(i)} + U_{(i+1)} + U_{(i+2)} \le U_{(i+3)}
  \] for all $1 \le i \le n-3$.

  We use the same method of rewriting the order statistics as sums of exponential variables, getting the inequality
  \begin{align*}
    \frac{X_1 + \dots + X_{i}}{S} + \frac{X_1 + \dots + X_{i+1}}{S} + \frac{X_1 + \dots + X_{i+2}}{S} &\le \frac{X_1 + \dots + X_{i+3}}{S} \\
    2(X_1 + \dots + X_i) + X_{i+1} &\le X_{i+3}
  \end{align*}
  which needs to be satisfied for all $i$, leaving $X_1, X_2$ and $X_3$ unbounded.

  Again, we integrate over all the variables to obtain the required probability, which we denote $P_n$. This gives us
  \begin{align*}
    P_n &= \int_{0}^{\infty} e^{-x_1}\int_{0}^{\infty} \dotsi \int_{2(x_1 + \dots + x_{n-3}) + x_{n-2}}^{\infty} e^{-x_n}\, dx_n\, \cdots\, dx_2\, dx_1, \\
        &= \int_{0}^{\infty} \int_{0}^{\infty} \dotsi \int_{2(x_1 + \dots + x_{n-3}) + x_{n-2}}^{\infty} e^{-(x_1 + \dots + x_n)}\, dx_n \, \cdots \, dx_2 \, dx_1.
  \end{align*}

  Let the integral obtained after $i$ integrations be denoted $I_i$. For example, we have \[
    I_1 = \int_{2(x_1+\dots+x_{n-3})+x_{n-2}}^{\infty} e^{-(x_1+\dots+x_n)}\, dx_n = e^{-(3(x_1+\dots+x_{n-3})+2x_{n-2}+x_{n-1})}
  .\]

  We now introduce three sequences, $(R_k)$, $(S_k)$ and $(T_k)$, to track the coefficients of the terms in the exponent during the inductive process of successive integrations. First, let $R_2=3$, $S_2=2$, $T_2=1$ so that the expression for $I_1$ above agrees with the $k=2$ case. We now derive recurrence relations from the coefficient updates in the exponent after each integration:
  \begin{align*}
    R_{k} &= R_{k-1} + 2T_{k-1}, \\
    S_k &= R_{k-1} + T_{k-1}, \\
    T_{k} &= S_{k-1}.
  \end{align*}
  This is seen inductively---the first integral $I_1$ evaluates to \[
    I_1 = e^{-(R_2(x_1 + \dots + x_{n-3}) + S_2 x_{n-2} + T_2 x_{n-1})},
  \] and we then assume that $I_i$ evaluates to
  \begin{equation} \label{eq3}
    I_i = \frac{1}{\prod_{k=1}^{i} T_k} e^{-(R_{i+1}(x_1 + \dots + x_{n-2-i}) + S_{i+1} x_{n-1-i} + T_{i+1}x_{n-i})}.
  \end{equation}
  From this inductive hypothesis, we obtain
  \begin{align*}
    I_{i+1} &= \int_{2(x_1+\dots+x_{n-3-i}) + x_{n-2-i}}^{\infty} I_i\, dx_{n-i} \\
            &= \frac{1}{\prod_{k=1}^{i+1} T_k} e^{-((R_{i+1} + 2T_{i+1})(x_1 + \dots + x_{n-3-i}) + (R_{i+1} + T_{i+1})x_{n-2-i} + S_{i+1}x_{n-1-i})},
  \end{align*}
  which closes the induction, showing that the recurrence relations hold. It remains to verify that $T_k$ indeed satisfies the Tribonacci recursion. Subtracting $S_k$ from $R_k$ gives us $R_k - S_k = T_{k-1}$, so $R_k = S_k + T_{k-1}$. From this, substitutions yield 
  \begin{align*}
    T_k = S_{k-1} &= R_{k-2} + T_{k-2} \\
                  &= (S_{k-2} + T_{k-3}) + T_{k-2} \\
                  &= T_{k-1} + T_{k-3} + T_{k-2},
  \end{align*}
  so our $T_k$ are exactly the Tribonacci numbers. Now, using \eqref{eq3}, we have \[
    I_{n-3} = \frac{1}{\prod_{k = 1}^{n-3} T_k} e^{-(R_{n-2}x_1 + S_{n-2}x_2 + T_{n-2}x_3)}
  ,\] leaving us with
  \begin{equation*}
    P_n = \int_{0}^{\infty} \int_{0}^{\infty} \int_{0}^{\infty} \frac{1}{\prod_{k = 1}^{n-3} T_k} e^{-(R_{n-2}x_1 + S_{n-2}x_2 + T_{n-2}x_3)} \, dx_3\, dx_2\, dx_1.
  \end{equation*}
  Evaluating this integral yields \[
    P_n = \frac{1}{T_1 \cdots T_{n-2} S_{n-2} R_{n-2}}
    .\] From the recursive definitions of $R$ and $S$ previously, we have that $S_{n-2} = T_{n-1}$, and that $R_{n-2} = T_{n} - T_{n-2}$, ultimately giving us \[
    P_n = \frac{1}{(T_n - T_{n-2})T_1\cdots T_{n-1}}
  ,\] as desired.
\end{proof}

\end{document}
